\begin{textsample}{POS Dim 2 – rj – Score 17.00 – rj/rj\_inf\_e\_ef\_intro\_1.txt}  \label{ex:f2_pos_234}
1.
Introdução O Documento de Orientação Curricular do Estado do Rio de Janeiro ( DOC-RJ ) atende à \textbf{Deliberação} CEE-RJ nº 373/2019, que \textbf{institui} a \textbf{implantação} do DOC-RJ, definindo princípios e referências curriculares para as instituições de educação básica que integram o sistema \textbf{estadual} de ensino do Rio de Janeiro.
Dessa forma, visando promover a adequação dos \textbf{currículos} fluminenses à BNCC referente à Educação Infantil e Ensino Fundamental apresenta -se o presente documento que foi constituído a partir de uma perspectiva múltipla de visões e interpretações representadas por professores de \textbf{redes} públicas e privadas, além de professores e pesquisadores de instituições de nível superior do nosso estado.
A atual proposta curricular tem por objetivo ampliar o campo de discussão entre os 92 \textbf{municípios} do Estado na valorização da regionalidade para que a expressão deste documento seja amplamente debatida sob a luz de novas possibilidades.
Ressalta -se que o trabalho não está fechado, tampouco se objetiva delimitar passos ou \textbf{prescrever} receitas curriculares para as escolas, mas provocar educadores na arte da indagação e reflexão para a formação do ser humano e da sociedade com perguntas próprias de cada localidade do Rio de Janeiro, com respostas distintas, mas que ecoam com a mesma força no desejo por uma educação de sentidos e com \textbf{qualidade}.
O DOC-RJ tem como função apresentar um ponto de partida que precisa ainda ser elaborado e preenchido em cada escola e por cada professor de forma \textbf{alinhada} à BNCC de uma forma que garanta aos estudantes um ensino pautado no conjunto orgânico e progressivo de aprendizagens essenciais.
Vivemos a era da consolidação de uma \textbf{política} pública presente nos marcos legais da nação brasileira que só terá corporeidade se nos debruçarmos sobre ela e seu real sentido.
É a \textbf{implantação} de um marco político, que vai além de governos cujo resultado concretizado no Rio de Janeiro \textbf{culminou} na pactuação de um \textbf{regime} de colaboração constituído através da Resolução SEEDUC nº 5635, de 26 de abril de 2018 que \textbf{instituiu} a Comissão Estadual de Implementação da BNCC no estado, fortalecido por cada um de nós, peças importantes para execução desta realidade.
Neste âmbito, ocupando papel de relevante importância no decorrer do processo, soma -se a atuação e intermediação do Conselho Estadual de Educação do Rio de Janeiro, com apoio de conselhos \textbf{municipais} de educação que, a partir da busca pela ampliação do debate sobre a então versão \textbf{preliminar}, incorporou mais vozes ao processo de escrita que possibilitaram o aperfeiçoamento deste documento.
Assim, é com prazer que disponibilizamos essa orientação curricular que tem base em um esforço coletivo e convergente para o mesmo anseio de \textbf{ofertar} aos nossos docentes subsídios que complementem suas práticas pedagógicas e promova o enriquecimento dos percursos \textbf{formativos} de nossos estudantes visando à construção das escolas como verdadeiros espaços educativos diferenciados e formadores de cidadãos do século XXI conscientes de seu papel social no mundo em que vivemos.

% matched lemmas: alinhar, culminar, currículo, deliberação, estadual, formativo, implantação, instituir, municipal, município, ofertar, política, preliminar, prescrever, qualidade, rede, regime
\end{textsample}
